\subsection{CHARACTERISTIC FUNCTIONS OF SETS}

Let E be a non-empty set and A a subset of E. The characteristic function of the subset A of E is the mapping \(\varphi_{A}\) of E into the set \(\{0, 1\}\) defined by

\[
\varphi_ {\mathbf {A}} (x) = 1 \quad \text { if } \quad x \in \mathbf {A}; \qquad \varphi_ {\mathbf {A}} (x) = 0 \quad \text { if } \quad x \in \mathbf {E} - \mathbf {A}.
\]

Clearly the relation \(\varphi_{\mathbf{A}} = \varphi_{\mathbf{B}}\) is equivalent to \(\mathbf{A} = \mathbf{B}\) . We have \(\varphi_{\mathbf{E}}(x) = 1\) for all \(x \in \mathbf{E}\) and \(\varphi_{\emptyset}(x) = 0\) for all \(x \in \mathbf{E}\) ; these are the only constant characteristic functions on \(\mathbf{E}\) . The following Proposition is an immediate consequence of the definitions:

PROPOSITION 7. For each pair of subsets A, B of a non-empty set E, we have

\begin{enumerate}
\def\labelenumi{(\arabic{enumi})}
\tightlist
\item
\end{enumerate}

\[
\varphi_ {\mathbf {E} - \mathbf {A}} (x) = 1 - \varphi_ {\mathbf {A}} (x),\tag{2}
\]

\[
\varphi_ {\mathbf {A} \cap \mathbf {B}} (x) = \varphi_ {\mathbf {A}} (x) \varphi_ {\mathbf {B}} (x),\tag{3}
\]

\[
\varphi_ {\mathbf {A} \cup \mathbf {B}} (x) + \varphi_ {\mathbf {A} \cap \mathbf {B}} (x) = \varphi_ {\mathbf {A}} (x) + \varphi_ {\mathbf {B}} (x)
\]

for all \(x \in E\) .

